% Page 054 — page imprimée 51
% Suite de la transcription de la lettre administrative reproduite dans le document.

\begin{enumerate}
\setcounter{enumi}{5}

\item
Le chef de groupe ARRIGNON, chef du groupe 2, ne quitte
pas un instant Monsieur ROTH quand celui-ci vient dans son groupe.
Il se rend compte en effet qu’il est dangereux de laisser le Pasteur seul
avec ses chefs d’atelier.

Monsieur ROTH dit en outre, à qui veut l’entendre, que le
gouvernement n’étant pas légitime, le serment que nous avons fait est
sans valeur.

\end{enumerate}

Nos jeunes chefs ont juste assez de leurs forces physiques,
intellectuelles, morales et spirituelles pour suivre et faire suivre
les directives bien précises qui doivent mener le pays vers son salut.

Dans ses dernières directives, B.O. 11\textsuperscript{o}, p. 5, le Général
Commissaire Général s’exprime ainsi : « \ldots\ un chef des chantiers n’est
pas un citoyen libre comme tout autre. C’est un fonctionnaire respon-
sable de l’éducation des jeunes. Il n’a pas le droit en conscience de
professer une opinion contraire aux directives du gouvernement. »

J’estime donc que Monsieur le Pasteur ROTH est un homme dange-
reux et je souhaite le voir partir.

Je vous prie, Monsieur le Pasteur, de bien vouloir agréer
l’expression de mes respectueuses salutations.

\begin{flushright}
signé : Louis HONORAT.
\end{flushright}
